\documentclass[12pt]{article}
\usepackage[utf8]{inputenc}
\usepackage[T1]{fontenc}
\usepackage{hyperref}
\title{Ash\'eninka Peren\'e Language Resources for K-12 Curriculum Development}
\author{Compiled by Claude Sonnet 5.5 (claude-sonnet-5-5) for Andr\'as Kornai}
\date{}
\begin{document}
\maketitle

\section{Introduction}
Ash\'eninka Peren\'e (ISO 639-3 prq; Arawakan; Peru) is an endangered language with roughly 1,000 speakers according to the survey brief. This search found scholarly and archival materials but no school textbooks, primers or curriculum documents. Nothing found indicates that the language is taught in schools. Note that the ELAR collection below is catalogued under a different code (arw), so check which variety a record covers.

\section{Primary Educational Resources}
No primers, textbooks or school materials were found in the search results. An ELDP project description mentions a planned bilingual Ash\'eninka Peren\'e--Spanish storybook, but no published edition was located.

\section{Linguistic and Reference Materials}
\begin{itemize}
\item Elena Mihas, \emph{A Grammar of Alto Perené (Arawak)} (Mouton Grammar Library 69, 2015), a comprehensive reference grammar. Front matter is available from James Cook University: \url{https://researchonline.jcu.edu.au/40717/}
\item Publisher page for the same grammar (access restricted): \url{https://www.degruyterbrill.com/document/doi/10.1515/9783110766301/html}
\item Berkeley SAPHON phonological inventory for the Perené dialect: \url{https://linguistics.berkeley.edu/saphon/en/inv/AsheninkaPe.html}
\end{itemize}

\section{Religious and Community Resources}
No scripture translations or community organisation pages for Perené were found in the search results.

\section{Digital and Audio Resources}
\begin{itemize}
\item ELDP/ELAR collection of audio and video recordings (2010) of myths, folk tales, oratory, songs, chants and riddles: \url{https://eldp.access.preservica.com/dk0146/}
\item NSF grant record for the documentation project (grammar, lexical database, recordings): \url{https://grantome.com/grant/NSF/BCS-0901196}
\end{itemize}

\section{Curriculum Development Context}
Existing material is documentary and descriptive. A curriculum would have to be built from the grammar and the recorded texts together with local speakers and educators. A trilingual dictionary was planned by the documentation project, but its publication was not confirmed.
\end{document}